\providecommand{\cL}{\mathscr{L}}
\providecommand{\cZ}{\mathcal{Z}}
\providecommand{\cH}{\mathcal{H}}
\providecommand{\fg}{\mathfrak{g}}
\providecommand{\cW}{\mathcal{W}}
\providecommand{\gSC}{\mathfrak{g}^{\mathrm{SC}}}
\providecommand{\gAmod}{\mathfrak{g}^{\mathrm{mod}}_\cA}
\providecommand{\Ydg}{Y^{\mathrm{dg}}}
\providecommand{\Zder}{\operatorname{Z}_{\mathrm{der}}}
\providecommand{\Abulk}{A_{\mathrm{bulk}}}
\providecommand{\Bbound}{B_{\partial}}
\providecommand{\mc}{\operatorname{MC}}
\providecommand{\Res}{\operatorname{Res}}
\providecommand{\RHom}{\operatorname{RHom}}

\chapter*{Preface}
\addcontentsline{toc}{chapter}{Preface}
\markboth{Preface}{Preface}

\vspace{1em}

\section*{I.\quad The bar coalgebra and its three functors}

\noindent
The bar complex $\barBch(\cA)$ on
$\FM_k(\C) \times \Conf_k^{<}(\R)$ is an $E_1$ coassociative
coalgebra: the differential encodes the chiral product via
residues on $\FM_k(\C)$, the deconcatenation coproduct records
the topological factorisation along $\R$. Three functors act on
this coalgebra and produce three distinct objects. Cobar
$\Omega(\barBch(\cA)) \simeq \cA$ inverts the construction.
Verdier duality
$\mathbb{D}_{\mathrm{Ran}}(\barBch(\cA)) \simeq \barBch(\cA^!)$
produces the line-side Koszul dual. Chiral Hochschild cochains
$C^\bullet_{\mathrm{ch}}(\cA,\cA)$, defined via the chiral
endomorphism operad $\End^{\mathrm{ch}}_\cA$ with spectral
parameters from $\FM_k(\C)$, compute the bulk observables. The
third functor is not a projection from $\barBch(\cA)$ in the
sense of the first two; it is a separate construction that takes
$\barBch(\cA)$ as input and produces a new object carrying its
own algebraic structure. Every chapter of this volume studies one
of these three outputs.

The $\mathrm{SC}^{\mathrm{ch,top}}$ structure does not live on
$\barBch(\cA)$. It lives on the \emph{pair}
$(C^\bullet_{\mathrm{ch}}(\cA,\cA),\, \cA)$: the chiral
Hochschild cochain complex acts on the boundary algebra $\cA$
via brace operations, and this bulk-to-boundary action is the
$\mathrm{SC}^{\mathrm{ch,top}}$ datum. A chiral algebra is a
point in the derived symplectic category: a Lagrangian
$\cL_\cA \hookrightarrow \cM_{\mathrm{vac}}(\cA)$. The bar
complex is the endomorphism coalgebra of this point; Koszul
duality is the symplectic complement; the genus tower is the
deformation of the point over the moduli of curves. Volume~I
computes the Taylor jets; this volume identifies the local
composition law.

The three functors play three genuinely different roles and never
coincide: cobar recovers the algebra up to quasi-isomorphism,
Verdier duality interchanges it with its line-side dual, and
chiral Hochschild cochains produce the bulk as a separate object
with its own module category. Conflating any two of them
collapses the architecture of the volume.

\section*{II.\quad Two directions}

Two operators $a(z_1, t_1)$ and $b(z_2, t_2)$ on $\C \times \R$
interact in two independent directions.

\smallskip
\noindent\textbf{The holomorphic direction} ($z_1 \to z_2$).
Collision produces the residue pairing governed by the
$\lambda$-bracket. Volume~I built the entire modular Koszul
theory from this mechanism: the bar differential
\[
d_{\barB}[s^{-1}a \,|\, s^{-1}b]
\;=\;
\Res_{z_1 = z_2}\!\bigl[
 a(z_1)\,b(z_2)\cdot\eta_{12}
\bigr],
\qquad
\eta_{12} = d\log(z_1 - z_2),
\]
the Arnold relation
$\eta_{12}\wedge\eta_{23} + \eta_{23}\wedge\eta_{31}
+ \eta_{31}\wedge\eta_{12} = 0$
forcing $D_\cA^2 = 0$ at genus~$0$, the curved lift
$d_{\mathrm{fib}}^2 = \kappa(\cA) \cdot \omega_g$ at higher genus,
and the universal Maurer--Cartan element
$\Theta_\cA = D_\cA - d_0$ whose projections are the five Vol~I
theorems.

\smallskip
\noindent\textbf{The topological direction.}
The operators separate: cut $\R$ at a point $t_*$ between $t_1$
and $t_2$. Observables on each half-line depend only on their
own operators. The \emph{deconcatenation coproduct} records this
separation:
\begin{equation}\label{eq:preface2-delta}
\Delta[s^{-1}a_1|\cdots|s^{-1}a_n]
\;=\;
\sum_{i=0}^{n}
[s^{-1}a_1|\cdots|s^{-1}a_i]
\;\otimes\;
[s^{-1}a_{i+1}|\cdots|s^{-1}a_n].
\end{equation}
The two directions are independent:
\begin{equation}\label{eq:preface2-coderivation}
\Delta \circ D_\cA
\;=\;
(D_\cA \otimes \id \;+\; \id \otimes D_\cA) \circ \Delta.
\end{equation}
The differential is a coderivation of the coproduct. Collisions
in $\C$ commute with cuts in $\R$. Together the two structures
exhibit the bar complex on
$\FM_k(\C) \times \Conf_k^{<}(\R)$: the operadic fingerprint of
the Swiss-cheese operad $\SCchtop$.

The two-structure thesis of this volume is the assertion that
these two pieces of data form a complete local description of a
three-dimensional holomorphic--topological theory. The
differential in~$\C$ is a first-order local fact about OPE
residues; the coproduct in~$\R$ is a zeroth-order groupoid fact
about concatenation. Neither is sufficient alone. The
coderivation compatibility between them is what pins the bar
complex to the Swiss-cheese operad, and every further structure
in this volume (the genus tower, the bulk-boundary-line triangle,
the spectral $R$-matrix, the PVA descent) projects from this
single compatibility identity.

\section*{III.\quad Three conceptual leaps}

Everything in both volumes is organised by three leaps away from
classical Koszul duality.

\smallskip
\noindent\textbf{(1) Commutative to $E_1/\Ainf$.}\enspace
Classical Koszul duality works for commutative/associative
algebras. The chiral setting replaces $\Com$ by the Swiss-cheese
operad $\SCchtop$: the holomorphic direction promotes commutative
multiplication to a tower of higher operations $m_k$ (the
$\Ainf$~Stasheff hierarchy on $\FM_k(\C)$), and the topological
direction promotes concatenation to the deconcatenation coproduct
on $\Conf_k^{<}(\R)$.

\smallskip
\noindent\textbf{(2) Genus~$0$ to modular.}\enspace
Classical Koszul duality is a genus-$0$ computation on the
formal disk: one differential, one cohomology. The modular
setting lifts the bar
complex over $\overline{\cM}_{g,n}$ and tracks the curved
differential $d_{\mathrm{fib}}^2 = \kappa(\cA) \cdot \omega_g$.
The genus tower is the deformation of the Lagrangian over the
moduli of curves.

\smallskip
\noindent\textbf{(3) Local to nonlocal.}\enspace
Chiral algebras in the sense of Beilinson--Drinfeld are local:
their $R$-matrix is derived from the local OPE. Genuinely
nonlocal ($E_1$) chiral algebras take the $R$-matrix as an
independent input: Yangians, Etingof--Kazhdan quantum vertex
algebras, dg-shifted factorisation quantum groups. The
distinction is provenance, not pole order.

\smallskip
These three leaps act independently on the classical picture,
and together they give the seven-part architecture of this
volume.

\section*{IV.\quad The geometric escalation}

Ten spaces carry bar chain models. Each gains one degree of
freedom over the last; each introduces a configuration space, an
algebraic structure, and a specific object on which that structure
lives. The volume is the ascent from Stage~$0$ to Stage~$9$.

\smallskip
\noindent\textbf{Stage~$0$: a point.}\enspace
The algebra $\cA$ at a single point. Classical Koszul duality
over $\operatorname{Spec}\mathbf{k}$: the operad is $\mathrm{Ass}$,
the structure is an associative product, and it lives on~$\cA$.

\smallskip
\noindent\textbf{Stage~$1$: the real line~$\R$.}\enspace
Configuration space $\Conf_n^{<}(\R)$ of $n$ ordered points.
The deconcatenation coproduct
$\Delta\colon T^c(s^{-1}\bar\cA) \to
T^c(s^{-1}\bar\cA) \otimes T^c(s^{-1}\bar\cA)$
makes $B(\cA)$ into a coassociative coalgebra over $\mathrm{Ass}^{\scriptstyle\text{\rm !`}}$.
The ordering lives here: every subsequent stage inherits it.

\smallskip
\noindent\textbf{Stage~$2$: the interval $[0,1]$ and the
half-line $\R_{\ge 0}$.}\enspace
Boundaries produce augmentations.
$B(\cA) = \mathbf{k} \otimes_\cA^{\mathbf{L}} \mathbf{k}$
is the two-sided bar complex on $[0,1]$ with trivial boundary
conditions; the one-sided bar $B(\cA, M)$ on $\R_{\ge 0}$
resolves a module $M$ at the endpoint. The algebraic structure
is the minimal $\Ainf$ resolution $W(\mathrm{Ass})$; it lives on
$B(\cA)$ and on the category of line operators
$\cC_{\mathrm{line}}$.

\smallskip
\noindent\textbf{Stage~$3$: the complex plane~$\C$.}\enspace
Configuration space $\FM_k(\C)$, the Fulton--MacPherson
compactification. OPE residues along the Arnold form
$\eta_{ij} = d\log(z_i - z_j)$ define the bar
differential~$D_\cA$. For $E_\infty$-chiral algebras (all
vertex algebras), the operad is the Beilinson--Drinfeld chiral
operad; for $E_1$-chiral algebras (Yangians, Etingof--Kazhdan
quantum vertex algebras), it is the ordered chiral associative
operad $\mathrm{Ass}^{\mathrm{ch}}$. The $\Ainf$ operations
$\{m_k\}$ are cogenerator projections of $D_\cA$; they live on
the bar differential of~$B(\cA)$.

\smallskip
\noindent\textbf{Stage~$4$: $\C \times \R$.}\enspace
The product of Stages $1$ and $3$. Configuration space
$\FM_k(\C) \times \Conf_k^{<}(\R)$. The differential $D_\cA$
from $\C$ and the coproduct $\Delta$ from $\R$ satisfy the
coderivation identity
$\Delta \circ D_\cA
= (D_\cA \otimes \id + \id \otimes D_\cA) \circ \Delta$:
collisions in $\C$ commute with cuts in $\R$.
The algebraic structure is an $E_1$ dg coassociative coalgebra
over $(\mathrm{Ass}^{\mathrm{ch}})^{\scriptstyle\text{\rm !`}}$
(the Koszul dual cooperad of the chiral associative operad);
it lives on $B(\cA)$.

\smallskip
\noindent\textbf{Stage~$5$: the upper half-plane~$\mathbb{H}$
($\partial\mathbb{H} = \R$).}\enspace
Two-coloured configuration spaces: bulk points in $\mathbb{H}$
collide via $\FM_k(\C)$; boundary points on
$\partial\mathbb{H} = \R$ collide via $\Conf_m^{<}(\R)$; bulk
points approach the boundary but not conversely. The operad is
$\SCchtop$ (Swiss-cheese, chiral--topological); the algebraic
structure is the $\SCchtop$ datum; it lives on the \emph{pair}
$(\cZ^{\mathrm{der}}_{\mathrm{ch}}(\cA),\, \cA)$ of bulk
derived center and boundary algebra.

\smallskip
\noindent\textbf{Stage~$6$: the formal disk~$D$.}\enspace
$\operatorname{Aut}(\cO)$-equivariance on the formal disk
defines a vertex algebra: the chiral endomorphism operad
$\End^{\mathrm{ch}}_\cA$ with its coordinate-independent
formulation. The algebraic structure is the local vertex algebra
datum; it lives on $\cA$ as a $D$-module on~$X$.

\smallskip
\noindent\textbf{Stage~$7$: annuli, higher-genus
surfaces~$\Sigma_g$.}\enspace
Sewing two boundary circles of an annulus produces the cyclic bar
complex; clutching raises genus. The bar complex lifts over
$\overline{\cM}_{g,n}$ and the fiberwise differential becomes
curved: $d_{\mathrm{fib}}^2 = \kappa(\cA) \cdot \omega_g$.
The algebraic structure is the partially modular extension of
$\SCchtop$; it lives on the universal Maurer--Cartan element
$\Theta_\cA$ and on the genus tower $\{F_g\}_{g \ge 1}$.

\smallskip
\noindent\textbf{Stage~$8$: the Drinfeld center.}\enspace
The chiral Hochschild cochain complex
$C^\bullet_{\mathrm{ch}}(\cA, \cA)$ carries an $E_2$-chiral
Gerstenhaber bracket from the chiral Deligne--Tamarkin principle:
the $\FM_k(\C)$ operations on the bar complex induce brace
operations on cochains, and the brace algebra is formal. The
algebraic structure is $E_2$-chiral; it lives on
$\cZ^{\mathrm{der}}_{\mathrm{ch}}(\cA)$.

\smallskip
\noindent\textbf{Stage~$9$: topologization.}\enspace
A conformal vector $T(z)$ at non-critical level provides a
Sugawara construction: $T(z) = \{Q, G(z)\}$, so
$\C$-translations become $Q$-exact, the complex structure on
$\C$ becomes invisible in cohomology, the two colours of
$\SCchtop$ collapse, and Dunn additivity
($E_2^{\mathrm{top}} \otimes E_1^{\mathrm{top}} = E_3^{\mathrm{top}}$)
promotes $\cZ^{\mathrm{der}}_{\mathrm{ch}}(\cA)$ to an
$E_3$-topological algebra independent of the complex structure
on~$\C$. Without conformal vector: stuck at $\SCchtop$.
At critical level $k = -h^\vee$: Sugawara undefined,
topologization fails.

\medskip
Stage~$9$ is the point of this volume.

\section*{V.\quad The bar complex as self-intersection}

The independence of the two directions is geometric. The bar
complex is a formal Koszul model of the Lagrangian
self-intersection $\cL_\cA \times_{\mathscr{M}}^h \cL_\cA$ in a
$(-2)$-shifted symplectic stack
(Theorem~\ref{thm:bar-is-self-intersection}). The holomorphic
direction is the Koszul differential of the Lagrangian
embedding; the topological direction is the groupoid diagonal;
the coderivation property is the compatibility of the two. Both
structures descend from the single fact that the bar complex
computes a derived fibre product.

The chiral algebra fixes a formal local bulk model: HKR for
Lagrangian embeddings recovers the shifted cotangent
$T^*[-1]\cL_\cA$ from boundary data; the actual ambient stack
requires the Darboux $1$-form
(Theorem~\ref{thm:holographic-reconstruction}).

\section*{VI.\quad The open primitive and the AP-OC distinction}

The bar complex is not the fundamental object. It is the
coalgebraic shadow of a more primitive datum: the open/closed
factorization dg-category~$\cC$ on the bordified curve
$\widetilde{X}_D$. The objects of $\cC$ are boundary
conditions; its morphisms are open-string states. A choice of
vacuum boundary condition $b$ produces a boundary algebra
$A_b = \End_\cC(b)$ as a \emph{chart}; the bar complex
$\barB(A_b)$ encodes how closed-sector operations twist the
open-sector composition. The Morita class of $\cC$, not the
quasi-isomorphism class of $A_b$, is the invariant.

Three objects are independently constructed from $\cC$ and must
never be conflated:
\begin{itemize}[nosep]
\item the \emph{bar complex} $\barB(A_b)$ classifies twisting
 morphisms (couplings between $A_b$ and its line-side Koszul
 dual $A^!_{\mathrm{line}}$);
\item the \emph{chiral derived center}
 $\cZ^{\mathrm{der}}_{\mathrm{ch}}(\cC)
 = \RHom_{\operatorname{Fun}(\cC,\cC)}(\operatorname{Id},
 \operatorname{Id})$ is the universal bulk;
\item the \emph{line-side Koszul dual algebra}
 $A^!_{\mathrm{line}}
 = D_{\operatorname{Ran}}(\barB(A_b))$, which models the line
 side on the chirally Koszul locus.
\end{itemize}

This is the \emph{AP-OC distinction}: the bar complex classifies
twisting morphisms; the bulk is the derived centre; the two live
on different levels of the categorical hierarchy. Conflating
them is the single most corrosive failure mode in the subject.
The slab model of the interface between bulk and boundary makes
this precise as a fibre functor
(Remark~\ref{rem:slab-fiber-functor} in
\S\ref{sec:ht-bulk-boundary-line}), and the interface
generalisation of Dimofte's construction
(Remark~\ref{rem:dimofte-interface-generalization} in the same
section) shows how the slab extends across changes of the
open/closed boundary data.

Modularity is not an additional axiom. It is trace plus
clutching on the open sector: the cyclic trace on $\End(b)$
seeds a Calabi--Yau structure; the annulus identification
$\int_{S^1} \cC \simeq \HH_\bullet(\cC)$ is the first modular
shadow; clutching raises genus. The shadow archetypes
G/L/C/M classify the $\Ainf$ complexity of the boundary algebra.

\section*{VII.\quad The Vol~I engine}

Volume~I is Koszul duality on a curve: the bar construction for
chiral algebras on an algebraic curve~$X$, with five main
theorems Theorems~A--D and~H of Vol~I. Theorems~A, B, and~C1
are unconditional on their stated loci; Theorem~H is on the
Koszul locus \textup{(}at generic affine
level in the non-abelian case\textup{)}; Theorem~D and the
scalar part C2 of Theorem~C
carry uniform-weight scope at $g \ge 2$, with a cross-channel
correction~$\delta F_g^{\mathrm{cross}}$ in the multi-weight
regime. Three structure theorems organise the
representation-theoretic content: \emph{algebraicity}
($H(t)^2 = t^4 Q_L(t)$, shadow generating function algebraic of
degree~$2$); \emph{formality identification}, where the shadow
obstruction tower agrees with the $L_\infty$ formality
obstruction tower at all degrees and the discriminant
$\Delta = 8\kappa S_4$ classifies algebras into four depth
classes~G/L/C/M; and \emph{complementarity}, proved under
perfectness and chain-level nondegeneracy, as a Lagrangian
splitting of factorisation homology.

Volume~II lifts this from one to three dimensions. The bar
differential $D_\cA$ is holomorphic factorisation on~$\C$; the
deconcatenation coproduct $\Delta$ is topological factorisation
on~$\R$; together they form an $E_1$ dg coassociative coalgebra on
$\FM_k(\C) \times \Conf_k^<(\R)$. The $\SCchtop$ structure
emerges in the chiral derived center: the pair
$(\cC^\bullet_{\mathrm{ch}}(\cA,\cA),\,\cA)$ carries
Swiss-cheese structure via brace operations. The modular
characteristic $\kappa(\cA)$ controls the curved $A_\infty$
deformation at genus $g \ge 1$; complementarity lifts to the
bulk-boundary-line triangle; the discriminant $\Delta$
classifies the complexity of bulk-boundary pairs. Classical
Koszul duality is genus~$0$ on the formal disk; Volume~I is
genus~$g$ on a curve; this volume is the dimensional lift to
$\C \times \R$.

\section*{VIII.\quad The Swiss-cheese operad $\SCchtop$}

The product $\FM_k(\C) \times \Conf_k^<(\R)$ is the space of
operations of a two-coloured operad:
\begin{alignat*}{2}
&\text{closed-to-closed:}\quad
&&\SCchtop(\mathsf{ch}^k;\,\mathsf{ch})
\;=\; \FM_k(\C), \\
&\text{mixed-to-open:}\quad
&&\SCchtop(\mathsf{ch}^k, \mathsf{top}^m;\,\mathsf{top})
\;=\; \FM_k(\C) \times E_1(m), \\
&\text{open-to-closed:}\quad
&&\SCchtop(\ldots, \mathsf{top}, \ldots;\,\mathsf{ch})
\;=\; \varnothing.
\end{alignat*}
Closed colour: holomorphic collisions on $\FM_k(\C)$. Open
colour: topological orderings on $E_1(m) = \Conf_m^<(\R)$.
Mixed: the product. Composition is FM substitution on the closed
factor, interval insertion on the open factor. The empty
open-to-closed component is directionality: bulk restricts to
boundary, not conversely.

The bar complex $(\barBch(\cA), D_\cA, \Delta)$ is an $E_1$ chiral
coassociative coalgebra (Theorem~\ref{thm:rosetta-e1-coalgebra}): the
differential encodes holomorphic factorization, the coproduct
encodes topological factorization. The $\SCchtop$ structure
emerges in the chiral derived center pair
$(C^\bullet_{\mathrm{ch}}(\cA,\cA),\, \cA)$.
The classical Swiss-cheese operad
is Koszul (Livernet); Kontsevich formality and bar-cobar transfer
upgrade this to a Quillen equivalence
$\Omega\mathbf{B}(\SCchtop) \xrightarrow{\sim} \SCchtop$
(Theorem~\ref{thm:homotopy-Koszul}). The Koszul dual cooperad
$(\SCchtop)^! = (\mathsf{Lie},\, \mathsf{Ass},\, \text{shuffle-mixed})$
(Proposition~\ref{thm:SC-self-duality}) exchanges the closed colour
$\mathsf{Com}$ with $\mathsf{Lie}$; the duality functor on
$\SCchtop$-algebras is an involution, and the open-colour dual
$\cA^!_{\mathrm{line}}$ carries the dual algebraic structure.

\section*{IX.\quad Three computations}

The machine is tested on three families of increasing complexity;
the full computations are carried out in the Introduction and in
Part~\ref{part:examples}.

\smallskip
\noindent\textbf{Heisenberg} (class~$\mathbf{G}$, formal).
One generator $J$, one OPE $J(z)J(w) \sim k/(z-w)^2$. The
$\lambda$-bracket is $\{J_\lambda J\} = k\lambda$, with
$m_2(J,J;\lambda) = k\lambda$ and all $m_{k\ge 3} = 0$. The
classical $r$-matrix is $r(z) = k/z$ (the $d\log$ kernel absorbs
one pole order from the double-pole OPE); the quantum $R$-matrix
is $R(z) = \exp(k\hbar/z)$. The chiral Koszul dual is
$\cA^!_{\mathrm{ch}} = \mathrm{Sym}^{\mathrm{ch}}(V^*)$ (note
$\cH_k^! \neq \cH_{-k}$); the open-colour dual is
$\cA^!_{\mathrm{line}} = Y(\mathfrak{u}(1))$; the line category
is Fock modules. Bulk: $\Abulk \simeq \Zder(\cH_k) = \cH_k$.
Genus~$1$: $\kappa(\cH_k) = k$, $F_1 = -k\log\eta(\tau)$. The
single scalar $k$ controls every entry.

\smallskip
\noindent\textbf{Kac--Moody} (class~$\mathbf{L}$, Lie-transverse).
$m_2(J^a, J^b;\lambda) = f^{ab}{}_c J^c + k\,\delta^{ab}\lambda$;
$m_3 \neq 0$ (Jacobiator homotopy); $m_{k\ge 4} = 0$. The
classical $r$-matrix is $r(z) = k\,\Omega/z$ (with the level $k$
surviving $d\log$ absorption, so that $r(z)$ vanishes at
$k = 0$); the KZ connection
$\nabla_{\mathrm{KZ}} = d - \frac{1}{k+h^\vee}\,\Omega\,\tfrac{dz}{z}$
is the flat connection on conformal blocks. Open-colour dual:
$\cA^!_{\mathrm{line}} = \Ydg(\fg)$. Genus~$1$:
$\kappa(\widehat{\fg}_k)
= \dim\fg\,(k + h^\vee)/(2h^\vee)$, with the Feigin--Frenkel
relation
$\kappa(\widehat{\fg}_k) + \kappa(\widehat{\fg}_{-k-2h^\vee})
= 0$.

\smallskip
\noindent\textbf{$\cW_3$} (class~$\mathbf{M}$, infinite tower).
The first family where the $\Ainf$ structure is genuinely
infinite: $m_2(T,T;\lambda)
= (c/12)\lambda^3 + 2T\lambda + \partial T$, and $m_3$, $m_4$
are both nonzero with a tower that does not terminate. The
cumulant coalgebra $\operatorname{Cum}_c(\cW_3)$ records the
first family whose kinematic complexity exceeds the Virasoro
baseline.

\section*{X.\quad The 3d MC element $\alpha_T$}

Volume~I has one MC element $\Theta_\cA$ in the modular
convolution algebra $\gAmod$. This volume has one MC element
$\alpha_T \in \mc(\gSC_T)$ in the Swiss-cheese convolution algebra
\[
\gSC_T
\;:=\;
\Hom_\Sigma\!\bigl(
 \mathbf{B}(\SCchtop),\;
 \End_{(\Bbound,\,\cC_{\mathrm{line}})}
\bigr).
\]
Its $L_\infty$-brackets are indexed by two-coloured stable trees;
the existence of $\alpha_T$ follows from homotopy-Koszulity of
$\SCchtop$ (Theorem~\ref{thm:homotopy-Koszul}), which gives
$\partial^2 = 0$ on $\mathbf{B}(\SCchtop)$.

Six projections of $\alpha_T$ organise the volume:
\begin{center}
\small
\renewcommand{\arraystretch}{1.15}
\begin{tabular}{lll}
\textbf{Projection} & \textbf{Structure} & \textbf{Section} \\
\hline
closed face
 & boundary $\Ainf$ algebra $\Bbound$
 & \S\ref{sec:axioms-Ainfty-chiral} \\[2pt]
open face
 & open-colour operations on $\cC_{\mathrm{line}}$
 & \S\ref{sec:line_operators} \\[2pt]
ordered $(2,0)$ collision
 & spectral $R$-matrix, YBE
 & \S\ref{sec:spectral_braiding} \\[2pt]
mixed $(1,1)$
 & bulk-to-boundary comparison
 & \S\ref{sec:ht-bulk-boundary-line} \\[2pt]
closed genus tower
 & $\Theta_\cA$ from Vol~I
 & \S\ref{subsec:rosetta-3d-mc} \\[2pt]
closed cohomology
 & PVA $\{-_\lambda -\}$
 & \S\ref{sec:PVA_descent}
\end{tabular}
\end{center}
The scoped identifications
$\cC_{\mathrm{line}} \simeq \cA^!_{\mathrm{line}}\text{-mod}$ and
$\Abulk \simeq \Zder(\Bbound)
\simeq \HH^\bullet_{\mathrm{ch}}(\cA^!_{\mathrm{line}})$
assemble these faces into the bulk-boundary-line triangle
(Theorems~\ref{thm:boundary-linear-bulk-boundary}
and~\ref{thm:lines_as_modules}) on the boundary-linear exact
sector and on the chirally Koszul locus respectively;
$\Abulk$ identifies with chiral Hochschild cochains of the
boundary algebra unconditionally for classes~$\mathbf{G}$,
$\mathbf{L}$, $\mathbf{C}$ and at genus~$0$ on class~$\mathbf{M}$
(Proposition~\ref{prop:uch-mellin-shadow}); on class~$\mathbf{M}$
at genus $g \geq 1$ the equivalence holds in the weight-completed
ambient via the MC5 pattern only, with chain-level closure in
$\mathrm{Ch}(\mathrm{Vect})$ remaining open as
Conjecture~\ref{conj:uch-gravity-chain-level}.

The colour decomposition
$\alpha_T = \alpha_T^{\mathrm{cl}} + \alpha_T^{\mathrm{mix}}
+ \alpha_T^{\mathrm{op}}$ splits the MC equation into three: a
pure closed equation (the Vol~I MC equation), a mixed equation
(the $r$-matrix and PVA bracket are $\Theta_\cA$-twisted
cocycles), and a pure open equation that closes the line
category.

\section*{XI.\quad PVA descent and the $\Ainf$ hierarchy}
\label{sec:preface-pva-descent}

The bar differential $D_\cA$ is a coderivation of the cofree
coalgebra $T^c(s^{-1}\bar\cA)$, hence determined by its
cogenerator projections $m_k \colon A^{\otimes k} \to
A((\lambda_1))\cdots((\lambda_{k-1}))$ of degree $2 - k$. The
identity $D_\cA^2 = 0$ expands into the $\Ainf$ Stasheff
relations; each relation is Stokes' theorem on $\FM_n(\C)$. The
higher operations are the unique solution of a tree-level
bootstrap equation (Theorem~\ref{thm:resolvent-principle}), with
$m_k$ a sum over $C_{k-1}$ planar binary trees.

On cohomology $H^\bullet(A, Q)$, the higher operations
$m_{k\ge 3}$ vanish by contractibility of $\Conf_k^{<}(\R)$. The
binary $m_2$ decomposes into a commutative product (regular
part) and a $\lambda$-bracket (singular part), and together they
form a $(-1)$-shifted Poisson vertex algebra
(Theorem~\ref{thm:cohomology_PVA}). The Jacobi identity comes
from three-face Stokes on $\FM_3(\C)$; skew-symmetry from the
exchange automorphism; Leibniz from regular-part factorisation.
The quantization problem inverts this descent: given a PVA,
construct an $\SCchtop$-algebra whose cohomological shadow
recovers it. Quillen equivalence of the bar-cobar adjunction
(Theorem~\ref{thm:homotopy-Koszul}) gives essential uniqueness of
the lift.

\section*{XII.\quad The bulk-boundary-line triangle}

A holomorphic--topological theory on $\C \times \R_{\ge 0}$
produces three algebraic objects: \emph{bulk} local operators
$\Abulk$ on the interior, \emph{boundary} local operators
$\Bbound$ on $\{0\} \times \C$, and \emph{line} operators
$\cC_{\mathrm{line}}$ along $\R_+ \times \{0\}$. On the
chirally Koszul locus,
\[
\Abulk
\;\simeq\;
\Zder(\Bbound),
\qquad
\cC_{\mathrm{line}}
\;\simeq\;
\cA^!_{\mathrm{line}}\text{-mod}.
\]

Two parallel line operators interact through the bulk, producing
a meromorphic braiding $R_{L_1, L_2}(z)$, and three lines give
the Yang--Baxter equation (Theorem~\ref{thm:YBE}) from the same
$\FM_3(\C)$ Stokes mechanism as PVA Jacobi. For Heisenberg,
$R(z) = e^{k\hbar/z}$. For Kac--Moody on evaluation modules,
$R(z)$ agrees with the quantum group $R$-matrix of
$U_q(\fg)$ at $q = e^{i\pi/(k+h^\vee)}$
(Theorem~\ref{thm:affine-monodromy-identification}).

\section*{XIII.\quad Curved genus expansion}

At genus $g \ge 1$ the fiberwise differential is curved:
$d^2_{\mathrm{fib}} = \kappa(\Bbound) \cdot \omega_g$ (Vol~I, Theorem~D). This
is not a defect but the content of the modular theory: the
nonvanishing of $d^2_{\mathrm{fib}}$ is how the genus tower enters the local
bar-cobar picture. The transgression algebra $B_\Theta$ absorbs
the curvature (adjoin a generator $\eta$ of degree~$1$ with
$d\eta = \Theta$), and the secondary anomaly $u := \eta^2$
controls the genus-$g$ Clifford completion
$\mathfrak{G}_g(B_\Theta)$. Away from $u = 0$ the higher genus
is Morita-invisible; on $u = 0$ the completion degenerates into
an exterior model for $H_1(\Sigma_g, \Z)$. The full development
is in Chapter~\ref{app:anomaly-completed-topological-holography}.

The modular bar differential $D = \sum_{g \ge 0} \hbar^g D^{(g)}$
is the Feynman transform of the modular operad
(Theorem~\ref{thm:formal-genus-expansion}); the canonical MC
element $\delta_\cA = D - D^{(0)}$ lives in the positive-genus
filtration; the universal obstruction $\mathrm{Ob}_g$ at
genus~$g$ and the variational modular class $[\varpi_g(\cA)]$
assemble into the modular PVA quantization theorem
(Theorem~\ref{thm:modular-bar}).

\section*{XIV.\quad The seven parts}

The seven parts of this volume trace the three conceptual leaps
of \S III through the Swiss-cheese machine.

\smallskip
\noindent\textbf{\emph{The Open Primitive}
(Part~\ref{part:swiss-cheese}).}\enspace
The open/closed factorization dg-category $\cC$, the Morita
viewpoint, the bar complex as the classifier of twisting
morphisms for a chosen boundary chart.

\smallskip
\noindent\textbf{\emph{The $E_1$ Core}
(Part~\ref{part:e1-core}).}\enspace
The $E_1$ bar coalgebra, the chiral derived center, the 3d MC
element $\alpha_T$, the ordered/associative $E_1$ side of Koszul
duality, and the dg-shifted factorization bridge to Yangians.

\smallskip
\noindent\textbf{\emph{Seven Faces of $r(z)$}
(Part~\ref{part:bbl-core}).}\enspace
Each projection of the degree-$(2,0)$ collision: spectral
braiding and YBE; DNP identification; the bulk-boundary-line
slab (\S\ref{sec:ht-bulk-boundary-line}, containing
Remark~\ref{rem:slab-fiber-functor} and
Remark~\ref{rem:dimofte-interface-generalization}); celestial
boundary transfer; affine half-space BV; planted-forest
synthesis on $\FM_3(X)$; and the Kontsevich integral.

\smallskip
\noindent\textbf{\emph{Characteristic Datum and Modularity}
(Part~\ref{part:examples}).}\enspace
Chiral Hochschild, brace structures, the relative Feynman
transform, modular PVA quantization, and the physical origins
of the holomorphic--topological twist.

\smallskip
\noindent\textbf{\emph{The Standard HT Landscape}
(Part~\ref{part:holography}).}\enspace
Yang--Mills boundary theory, instanton screening, celestial
holography, logarithmic HT monodromy, anomaly completion, the
holographic reconstruction theorem, and the modular bootstrap.

\smallskip
\noindent\textbf{\emph{Three-Dimensional Quantum Gravity}
(Part~\ref{part:gravity}).}\enspace
The climax of the volume: $SL(2,\R) \times SL(2,\R)$ Chern--Simons
with $\Bbound = \mathrm{Vir}_{6k}$ is the Virasoro archetype of
class~$\mathbf{M}$, and its gravitational interpretation supplies
the perturbative finiteness, soft graviton theorems, and
symplectic polarization statements. The chiral dual at $c = 13$
(equivalently $k = 13/6$) is the structural antipode to the
matter--ghost critical dimension $c = 26$; the $\Ainf$ structure
is infinite-depth, the mixed archetype of Vol~I.

\smallskip
\noindent\textbf{\emph{The Frontier}
(Part~\ref{part:frontier}).}\enspace
Open conjectures: completeness of the higher-genus Lagrangian,
the global geometry of $\cM_{\mathrm{vac}}(\cA)$, the
Calabi--Yau origin story of Volume~III, and the single open
problem of the programme.

\section*{XV.\quad Signposts to the main theorems}

The principal results of this volume, each stated once in the
body and referenced throughout:
\begin{itemize}[leftmargin=2em,itemsep=2pt]
\item $E_1$ coalgebra structure on the bar complex and the
 chiral derived center as $\SCchtop$ datum
 (\S\ref{sec:rosetta-stone}).
\item Homotopy-Koszulity of $\SCchtop$ and the Quillen
 equivalence of bar and cobar
 (Theorem~\ref{thm:homotopy-Koszul}).
\item PVA descent of $m_2$ to a $(-1)$-shifted Poisson vertex
 algebra (Theorem~\ref{thm:cohomology_PVA}).
\item Spectral Yang--Baxter equation
 (Theorem~\ref{thm:YBE}).
\item Bulk-boundary-line triangle
 (Theorems~\ref{thm:boundary-linear-bulk-boundary}
 and~\ref{thm:lines_as_modules}).
\item Existence and universality of the 3d MC element $\alpha_T$
 (Theorem~\ref{thm:3d-universal-mc}).
\item Affine monodromy identification with $U_q(\fg)$
 (Theorem~\ref{thm:affine-monodromy-identification}).
\item Bar as self-intersection and holographic reconstruction
 (Theorems~\ref{thm:bar-is-self-intersection}
 and~\ref{thm:holographic-reconstruction}).
\item Complete spectral Drinfeld strictification
 (Theorem~\ref{thm:complete-strictification}).
\item General half-space BV quantization and the doubling
 theorem (Theorems~\ref{thm:general-half-space-bv}
 and~\ref{thm:doubling-rwi}).
\item Bulk-boundary-line factorization
 (Theorem~\ref{thm:bulk-boundary-line-factorization}).
\item Modular PVA quantization and the Feynman transform of the
 modular operad (Theorems~\ref{thm:modular-bar}
 and~\ref{thm:formal-genus-expansion}).
\item Logarithmic HT monodromy synthesis
 (Theorem~\ref{thm:synthesis}).
\end{itemize}
Each signpost is proved in a single chapter; the reader may
follow them in any order, because every chapter is a different
projection of the single coalgebra $\barBch(\cA)$.

\section*{XVI.\quad The local geometry of a point}

A chiral algebra is a point in the derived symplectic category:
a Lagrangian $\cL_\cA \hookrightarrow \cM_{\mathrm{vac}}(\cA)$.
The bar complex is the endomorphism coalgebra of this point,
Koszul duality is the symplectic complement, the genus tower is
the deformation of the point over the moduli of curves, and the
$\mathrm{SC}^{\mathrm{ch,top}}$ structure on the pair
$(C^\bullet_{\mathrm{ch}}(\cA,\cA),\, \cA)$ is the local
composition law governing bulk-to-boundary action.
The entire content of both volumes is the local geometry of the
derived symplectic category at this single point. Volume~I
computes the Taylor jets of the Lagrangian embedding. This
volume identifies the local composition law. What remains is
the global geometry of the stack: the completeness question for
the higher-genus Lagrangian, the moduli of Lagrangians, and the
higher-categorical descent. Volume~III asks how this local
geometry arises from Calabi--Yau categories.
